% ===== PRÉAMBULE CANONIQUE — à coller VERBATIM en tête de CHAQUE .tex =====
\documentclass[11pt,a4paper]{article}
\usepackage[utf8]{inputenc}\usepackage[T1]{fontenc}\usepackage{lmodern}
\usepackage[french]{babel}
\usepackage{amsmath,amssymb,mathtools}\usepackage{icomma}
\usepackage[version=4]{mhchem}          % \ce{...} : équations & formules
\usepackage{siunitx}\sisetup{locale=FR} % \qty{}{}, \si{}, \ang{}
\usepackage{chemfig}                    % \chemfig{...} : structures développées
\usepackage{xcolor}\usepackage{colortbl}
\usepackage{tikz}\usepackage{pgfplots}\pgfplotsset{compat=1.18}
\usepackage[most]{tcolorbox}\usepackage{enumitem}\usepackage{array,booktabs}
\usepackage[a4paper,margin=2.2cm]{geometry}
\usetikzlibrary{arrows.meta,calc,angles,quotes,positioning,patterns,backgrounds,shapes.geometric}
\definecolor{primary}{RGB}{222,49,99}\definecolor{primarydark}{RGB}{150,22,55}
\definecolor{defcol}{RGB}{0,114,178}\definecolor{methcol}{RGB}{52,92,148}
\definecolor{excol}{RGB}{0,135,90}\definecolor{attncol}{RGB}{200,40,55}
\tcbset{boxbase/.style={enhanced,breakable,boxrule=0.9pt,arc=2.5pt,
  left=3mm,right=3mm,top=2.3mm,bottom=2.3mm,fonttitle=\bfseries,coltitle=white}}
% --- boîtes TUTORIEL (noms EXACTS, ne pas renommer) ---
\newtcolorbox{cle}[1][]{boxbase,colback=defcol!6,colframe=defcol,
  title={\textsc{À retenir}\ifx\relax#1\relax\relax\else\textmd{ : \itshape #1}\fi}}
\newtcolorbox{methode}[1][]{boxbase,colback=methcol!7,colframe=methcol,sharp corners,
  title={\textsc{Méthode}\ifx\relax#1\relax\relax\else\textmd{ --- \itshape #1}\fi}}
\newtcolorbox{exemplebox}[1][]{boxbase,colback=excol!5,colframe=excol,
  title={\textsc{Exemple}\ifx\relax#1\relax\relax\else\textmd{ : \itshape #1}\fi}}
\newtcolorbox{piege}[1][]{boxbase,colback=attncol!6,colframe=attncol,
  title={\textsc{Piège}\ifx\relax#1\relax\relax\else\textmd{ : \itshape #1}\fi}}
\newtcolorbox{remarque}[1][]{boxbase,colback=black!4,colframe=black!45,
  title={\textsc{Remarque}\ifx\relax#1\relax\relax\else\textmd{ : \itshape #1}\fi}}
% --- boîtes EXERCICE / CORRIGÉ (noms EXACTS, auto-comptées) ---
\newtcolorbox[auto counter]{exo}[1][]{boxbase,colback=primary!4,colframe=primary,
  title={\textsc{Exercice}~\thetcbcounter\ifx\relax#1\relax\relax\else\textmd{ --- \itshape #1}\fi}}
\newtcolorbox[auto counter]{corrige}[1][]{boxbase,colback=excol!4,colframe=excol,
  title={\textsc{Corrigé}~\thetcbcounter\ifx\relax#1\relax\relax\else\textmd{ --- \itshape #1}\fi}}
\newcommand{\R}{\mathbb{R}}\newcommand{\N}{\mathbb{N}}\newcommand{\Z}{\mathbb{Z}}\newcommand{\ds}{\displaystyle}
% (un \usepackage{fancyhdr}+en-tête est possible ; il est ignoré côté web)
% =========================================================================
\begin{document}

\begin{corrige}[nommer un hydroxyde (base)]
\begin{enumerate}[label=\alph*)]
  \item \ce{NaOH} : hydroxyde de sodium.
  \item \ce{Ca(OH)2} : hydroxyde de calcium.
  \item \ce{Al(OH)3} : hydroxyde d'aluminium.
\end{enumerate}
\end{corrige}

\begin{corrige}[écrire un hydroxyde]
Le nombre de groupes \ce{OH} égale la valence du métal.
\begin{enumerate}[label=\alph*)]
  \item hydroxyde de potassium : \ce{K+} $\Rightarrow$ \ce{KOH}.
  \item hydroxyde de magnésium : \ce{Mg^2+} $\Rightarrow$ \ce{Mg(OH)2}.
  \item hydroxyde de zinc : \ce{Zn^2+} $\Rightarrow$ \ce{Zn(OH)2}.
\end{enumerate}
\end{corrige}

\begin{corrige}[nommer un sel neutre binaire]
Anion (-ure) d'abord, puis le cation.
\begin{enumerate}[label=\alph*)]
  \item \ce{NaCl} : chlorure de sodium.
  \item \ce{K2S} : sulfure de potassium.
  \item \ce{CaCl2} : chlorure de calcium.
\end{enumerate}
\end{corrige}

\begin{corrige}[écrire un sel neutre binaire]
\begin{enumerate}[label=\alph*)]
  \item chlorure de sodium : \ce{Na+} + \ce{Cl-} $\Rightarrow$ \ce{NaCl}.
  \item sulfure de sodium : \ce{Na+} + \ce{S^2-} $\Rightarrow$ \ce{Na2S}.
  \item bromure de potassium : \ce{K+} + \ce{Br-} $\Rightarrow$ \ce{KBr}.
\end{enumerate}
\end{corrige}

\begin{corrige}[les hydrures]
L'ion hydrure \ce{H-} est monovalent.
\begin{enumerate}[label=\alph*)]
  \item \ce{LiH} : hydrure de lithium.
  \item \ce{NaH} : hydrure de sodium.
  \item hydrure de calcium : \ce{Ca^2+} + \ce{H-} $\Rightarrow$ \ce{CaH2}.
\end{enumerate}
\end{corrige}

\end{document}
